\section{An infinite family attaining the exact maximum}\label{sec:family}
\begin{theorem}\label{thm:family}
For every integer $m\ge1$, take
\[
 \Delta_m=-2^{\,6\cdot7^{m-1}},\qquad h_m=2^{\,3\cdot7^{m-1}-2}.
\]
Then $h(\Delta_m)=h_m$ and $s(\Delta_m)=h_m+1$. This statement is independent of GRH.
\end{theorem}
The family in \cite{sparseparallel} has discriminants $-2^{2+6\cdot7^{m-1}}$ for $m\ge2$. The discriminants here differ by a factor of four. In particular, every choice of $h_m$ distinct nonnegative powers of such a $j$ is a basis for $\Q(j)$ over $\Q$. Dependence would give a polynomial with at most $h_m$ terms. Taking the exponents $0,1,\ldots,h_m$ shows that the minimum is attained. At $m=2$ it is $524{,}289$. At $m=3$ it is $2^{145}+1$.

\begin{proof}
Write $M=7^m$, $h=2^{3\cdot7^{m-1}-2}$, $f=2h$. The order of conductor $f$ in $\Q(i)$ has discriminant $-4f^2=\Delta_m$. Its class number is $f/2=h$, by the order class number formula: the Gaussian field has class number one, the unit index is two, and 2 ramifies.

The prime 7 is inert in $\Q(i)$ and coprime to $f$. The ring class field is unramified above 7. By Deuring reduction every conjugate of $j$ reduces to a supersingular invariant. In characteristic 7 the only such invariant is $1728\equiv6$. To check uniqueness directly, the Hasse invariant of $y^2=x^3+Ax+B$ in characteristic 7 is the coefficient of $x^6$ in $(x^3+Ax+B)^3$, namely $3B$. Supersingularity forces $B=0$. Thus all $h$ distinct conjugates lie in the same disk $6(1+7\mathcal O)$ in a common unramified completion.\cite{cmred}

Put $w(n)=n-v_7(n!)$. We show
\[
  w(i)<w(h)\ (0\le i<h),\qquad w(h)<w(i)\ (i>h).
\]
Because $2^3=1+7$, raising this equality to the power $7^{m-1}$ gives $4h\equiv1\pmod M$. For an interval of $r$ consecutive positive integers adjoining the cut, either $[h-r+1,h]$ (when $r\le h$) or $[h+1,h+r]$, it suffices to prove that the sum of the 7-adic valuations in the interval is strictly less than $r$.

First take $1\le r<M/5$. Its elements have the form $h+j$ with $|j|\le r$. The nonzero integer $1+4j$ has $|1+4j|\le4r+1<M$, so $v_7(h+j)=v_7(1+4j)$. In either arithmetic progression of $r$ consecutive choices of $j$, there are at most $r/7^a+1$ multiples of $7^a$ among the values $1+4j$. Summing over $a\le\log_7(4r+1)$ bounds the valuation sum by
\[
    r/6+\log_7(4r+1)<r .
\]
Indeed $4r+1\le5^r$ and $5^6<7^5$.

If $r\ge M/5$, direct counting of multiples in the original interval instead gives $r/6+\log_7(h+r)$. When $r\le h$, we have
\[
 \log_7(h+r)\le\log_7(2h)<M/6\le5r/6 ,
\]
where $2^{18}<7^7$ implies $(3/7)\log_7 2<1/6$. When $r>h$, use $h+r<2r\le 5^r$ and $\log_7 5<5/6$. These exhaust the two intervals and prove the strict inequalities for $w$.

It follows that $\min_{n\ge h} w(n)>\max_{0\le i<h}w(i)$. Lemma~\ref{lem:disk} prevents a polynomial with at most $h$ terms from vanishing at $h$ points of the disk. The degree-$h$ class polynomial supplies the opposite inequality.
\end{proof}

\paragraph{How frequently a single residue helps.}
For any order with $(\Delta/7)=-1$ or $(\Delta/13)=-1$, the same argument through Corollary~\ref{cor:coarse} yields
$s(\Delta)\ge h-O(\log(2h))$, independently of GRH. In characteristic 13 there is also exactly one supersingular invariant, and it is nonzero (the standard count is $\lfloor p/12\rfloor$ plus the corrections for 0 and 1728). The conductor is prime to an inert prime detected by the symbol. Among order discriminants counted once by absolute value, the stated union has natural density $1-(1-3/7)(1-6/13)=9/13$ by the Chinese remainder theorem. This is the density of the sufficient congruence conditions.
\subsection*{Full support for a positive proportion of prime conductors}
Let $\ell\ne2,7$ be prime, $\Delta=-4\ell^2$ and
$h=(\ell-\chi_{-4}(\ell))/2$. All $h$ conjugates reduce to the same
nonzero supersingular $j$ in characteristic $7$. Write
$w(n)=n-v_7(n!)$. By the local argument, $s(\Delta)=h+1$ whenever
\[
 w(u)<w(h)<w(v)\qquad (0\le u<h<v).
\]
Among primes $\ell$, this sufficient condition has a relative natural
density $d>0.6133$. Consequently the lower relative density of primes
giving $s(-4\ell^2)=h+1$ exceeds $61.33\%$. This result is independent
of Dirichlet GRH.

\smallskip\noindent\emph{Proof.}
An obstruction at distance $r$ requires
$\sum_{t\in I}v_7(t)\ge r$ on at least one of
$I=[h-r+1,h]$ (when $r\le h$) and $I=[h+1,h+r]$.
Let $b=\max_{t\in I}v_7(t)$. Counting multiples of $7^a$ in the
interval gives
\[
 \sum_{t\in I}v_7(t)\le\sum_{a=1}^b(r/7^a+1)\le r/6+b.
\]
Hence $q_r=7^{\lceil5r/6\rceil}$ divides $h+j$ for an offset
$1-r\le j\le r$. For fixed $r$, the prime number theorem in
progressions, applied separately to each sign $\sigma=\chi_{-4}(\ell)$
with $\ell=2h+\sigma$, bounds the relative density of the obstruction by
$u_r=2r/\varphi(q_r)$. Indeed there are at most $2r$ admissible classes
of $\ell$ per sign modulo $4q_r$, each of density $1/(2\varphi(q_r))$.

For completeness, these fixed-modulus bounds can be summed to control
arbitrarily distant obstructions. For primes $\ell\le X$ and
$q_r\le\sqrt X$, Brun--Titchmarsh bounds the count for each coprime
residue modulo $q_r$ by $2X/(\varphi(q_r)\log(X/q_r))$.
Ignoring the congruence modulo $4$, both signs and all offsets cost at
most $4r$ such classes. Thus for $r>R$ in this range the total is at
most $(X/\log X)\sum_{r>R}16r/\varphi(q_r)$. If $q_r>\sqrt X$,
a positive integer $h+j$ in the interval is divisible by
$7^{\lfloor\log_7\sqrt X\rfloor}$. Necessarily
$5r/6\le\log_7(h+r)$. For $r>h$ this fails, since $2r\le2^r$ and
$\log_7 2<5/6$. For $r\le h$ it implies $r=O(\log X)$.
There are $O(\sqrt X)$ possible multiples and $O(\log X)$ offsets,
hence $O(\sqrt X\log X)$ possible primes. This is $o(\pi(X))$.
The sieve inequality used here is due to Montgomery and Vaughan,
\emph{The large sieve}, Mathematika 20 (1973), 119--134
(\url{https://doi.org/10.1112/S0025579300004708}).

For fixed $R$, the event of no obstruction at distances $r\le R$ has
a density $d_R$, computable using complete residue classes modulo
$4\cdot7^R$: every valuation may be truncated at $R$ without affecting
a comparison with $r\le R$. The finitely many $h<R$ are immaterial.
The tail estimate just proved and convergence of $\sum u_r$ imply that
the density of no obstruction at any distance exists and equals
$d=\lim_Rd_R$. Applying the fixed-$r$ union bound first on
$R<r\le T$ and then letting $T\to\infty$ improves the error to
\[
                  0\le d_R-d\le\sum_{r>R}u_r.
\]
For instance, if $a=R+1$ and $b=(1+6/a)7^{-5}$, then
$u_{r+6}/u_r=(1+6/r)7^{-5}$ gives the explicit bound
$\sum_{r>R}u_r\le (u_a+\cdots+u_{a+5})/(1-b)$.

The supplied exact calculation gives
\[
 d_{10}=\frac{14143023}{23059204},\qquad
 d_{40}=\frac{33191837301111812268560181397}
              {54116956037952111668959660849},\qquad
 0\le d_{40}-d<5.6\cdot10^{-28}.
\]
To reproduce it, start with the twelve pairs consisting of a sign
$\sigma$ and a residue $h$ modulo $7$ with $7\nmid 2h+\sigma$.
Each pair has mass $1/12$. At depth $v$, compute each interval sum
with valuations capped at $v$. Discard the class if a sum already
reaches its interval length. Keep the class if every sum is smaller
and every valuation below $v$. Otherwise split into its seven lifts,
each with one seventh of the mass. At depth $R$, every ambiguous
interval already has truncated sum at least $R$ and is discarded.
Summing the kept masses yields $d_R$.
Thus $d$ rounds to $0.6133352599845868$ at sixteen places after the
decimal point. The exact-support set could contain further primes
for which this sufficient condition fails.\hfill$\square$

